\section{Consistencia macroscópica}\label{sec:macro}

Supongamos que todavía queda una fracción fija de las filas: $k\ge\delta_0 n$.
En la clase cuantitativa de la Sección~\ref{sec:B3T05}, queremos comparar dos operaciones.
Una consiste en fijar la reina en $c$ y volver a calcular el predictor en el hijo.
La otra consiste en variar la actividad de $c$ en el padre y tomar la derivada del predictor.
La segunda operación mantiene el soporte, mientras que la primera elimina casillas.

Las Secciones~\ref{sec:B4T01} a~\ref{sec:B4T04} tratan el cambio de soporte y construyen
el punto estacionario del hijo. Las Secciones~\ref{sec:B4T06} a~\ref{sec:B4T09} calculan los
términos de orden $1/k$ de la comparación y muestran cómo se cancelan en el resto.
La Sección~\ref{sec:B4T05} obtiene así la cota para cada hijo y, al sumar sobre la fila,
la cota para $\log\rho(B)$. Los pasos intermedios se reúnen en los Lemas~\ref{lem:hijo},
\ref{lem:escalar}, \ref{lem:prefactor} y~\ref{lem:modocero}, numerados por subsección.
